% Teaching manuscript. Edit this TeX source; Markdown is compatibility navigation only.

\StudyChapter{CS5222-Chapter03Part3}{Chapter 3 (Part 3) · TCP 可靠传输、流量控制、连接管理与拥塞控制}
\StudyTarget{CS5222-Chapter03Part3__root-section}

\StudyRepoRef{CS5222/course-notes/README.md}{}{课程目录} · \StudyRef{CS5222-Chapter03}{root}{上一篇：Chapter 3 Part 1–2} · \StudyRef{CS5222-Chapter04}{root}{下一章：Chapter 4 网络层数据平面} · \StudyRef{CS5222-Tutorial06}{root}{Tutorial 6} · \StudyRef{CS5222-Tutorial07}{root}{Tutorial 7}

\StudyHeading{1}{导读｜把“不撑满接收端缓存”与“不压垮网络路由器”分开}
\StudyTarget{CS5222-Chapter03Part3__overview}

在上一篇\StudyRef{CS5222-Chapter03}{root}{Chapter 3 Part 1–2}里，我们已经把传输层的编号方式从“按包计数（GBN/SR）”切换到了 TCP 的“按字节计数”，并用指数加权移动平均算出了超时间隔（RTO）。然而要让一条真实的跨洲 TCP 连接平稳运转，还必须回答四个紧接着出现的工程问题：第一，报文段丢失时怎样用单一计时器和冗余确认尽快恢复；第二，发送方发得太快、接收方应用程序来不及读缓存时怎样减速（流量控制）；第三，通信前后怎样可靠地建立和拆除状态（连接管理）；第四，成千上万条连接共享瓶颈路由器时怎样探测可用带宽并公平退让（拥塞控制）。

\textbf{基础补课（可跳过）：}对字节区间首尾计数或流水线时间轴不熟时，先回看\StudyRef{Foundations-NetworkBasics}{messages}{报文与字节区间}和\StudyRef{Foundations-NetworkBasics}{timeline}{时间轴与流水线}。

\begingroup\small\setlength{\tabcolsep}{3pt}\renewcommand{\arraystretch}{1.22}
\begin{longtable}{@{}>{\raggedright\arraybackslash}p{\dimexpr0.26\linewidth-1.33333\tabcolsep\relax}>{\raggedright\arraybackslash}p{\dimexpr0.48\linewidth-1.33333\tabcolsep\relax}>{\raggedright\arraybackslash}p{\dimexpr0.26\linewidth-1.33333\tabcolsep\relax}@{}}
\toprule
模块主题 & 核心机制与数量关系 & 配套练习与真题 \\
\midrule\endfirsthead
模块主题 & 核心机制与数量关系 & 配套练习与真题 \\
\midrule\endhead
\bottomrule\endfoot
TCP 重传与快速重传 & 单一定时器绑定最早未确认字节 \texttt{SendBase}；收到 \textbf{3 个冗余 ACK}（共 4 次相同确认号）立即重传缺失段 & \StudyRef{CS5222-Tutorial06}{q2}{T6 Q2}、\StudyRef{CS5222-Tutorial08}{p1}{T8 P1} \\
流量控制（\texttt{rwnd}） & \texttt{LastByteSent - LastByteAcked <= min(cwnd, rwnd)}；保护\textbf{接收端主机缓存} & \StudyRef{CS5222-Tutorial08}{p1}{T8 P1d–f} \\
三次握手与四次挥手 & \texttt{SYN / SYN+ACK / ACK}；双向独立 \texttt{FIN} 与主动关闭端 \texttt{TIME\_WAIT}（$2\times\mathrm{MSL}$） & 本篇 §3 \\
TCP Reno 与 Tahoe & 3 dup ACKs：Reno 置 $\mathrm{ssthresh}=\lfloor \mathrm{cwnd}/2\rfloor,\; \mathrm{cwnd}=\mathrm{ssthresh}+3$；Tahoe 置 $\mathrm{cwnd}=1$。Timeout：二者均置 $\mathrm{cwnd}=1$ & \StudyRef{CS5222-Tutorial07}{p1}{T7 P1}、2021A Q4 \\
稳态吞吐率与公平性 & 锯齿稳态平均吞吐率 $\approx \frac{3}{4}\frac{W}{\mathrm{RTT}}$；AIMD 使同 RTT 单连接收敛至等分，但并行连接与 UDP 可打破公平 & \StudyRef{CS5222-Tutorial07}{p2}{T7 P2} \\
\end{longtable}\endgroup


\StudyHeading{1}{1. TCP 可靠数据传输与快速重传（Fast Retransmit）}
\StudyTarget{CS5222-Chapter03Part3__tcp-rdt}

下层 IP 网络只提供尽力而为（best-effort）服务：数据报可能丢失、损坏或乱序到达。TCP 在此之上提供按序、无缺口的可靠字节流，依靠三项机制协同工作：\textbf{流水线发送（pipelining）}、\textbf{累积确认（cumulative acknowledgements）}与\textbf{绑定最早未确认字节 \StudyInline{SendBase} 的单一重传定时器}。

\begin{StudyBlock}{简化 TCP 发送端的三类驱动事件}{Chapter 3 Part 3 第 3–5 张幻灯片}
\begin{enumerate}
\item \textbf{从上层收到待发数据}：封装起始字节序号为 \StudyInline{NextSeqNum} 的报文段交付 IP，令 $\mathrm{NextSeqNum} \leftarrow \mathrm{NextSeqNum} + L$（$L$ 为本段数据字节数）；若重传定时器当前未运行，则启动定时器。
\item \textbf{定时器超时（Timeout）}：仅重传\textbf{当前序号最小的那一个未确认报文段}（这与 GBN 超时把整个窗口全部重传显著不同），随后重启定时器。
\item \textbf{收到确认号为 $y$ 的 ACK}：由于采用累积确认，$y$ 表示接收方已连续收妥 $y-1$ 及之前的所有字节。若 $y > \mathrm{SendBase}$，说明确认前沿向前推进，更新 $\mathrm{SendBase} \leftarrow y$；此时若仍有已发未确认字节在途，则重启定时器以守护新的 \StudyInline{SendBase}，否则停止定时器。
\end{enumerate}
\end{StudyBlock}

课件第 6–9 张幻灯片依次展示了四类典型收发时序，作图或阅读题干时需先看清丢的是数据段还是确认段：
\begin{itemize}
\item \textbf{情景一：确认段丢失（Lost ACK）}。主机 A 发送 \StudyInline{Seq=92, 8 bytes}，主机 B 收妥后发回 \StudyInline{ACK=100}，但该确认在途中丢失。A 的定时器到期后重发 \StudyInline{Seq=92, 8 bytes}；B 对比序号发现字节 $92..99$ 已在接收缓存中，于是丢弃重复副本，并再次发回 \StudyInline{ACK=100}。
\item \textbf{情景二：过早超时（Premature Timeout）}。A 背靠背发出 \StudyInline{Seq=92, 8 bytes} 与 \StudyInline{Seq=100, 20 bytes}，两段都顺利到达 B，B 依次发回 \StudyInline{ACK=100} 和 \StudyInline{ACK=120}，只是回程稍慢。A 在收到 \StudyInline{ACK=100} 之前超时重传了 \StudyInline{Seq=92, 8 bytes}；当该重传段抵达 B 时，B 早已收齐至 $119$ 号字节，因此再次回传当前最新期望值 \StudyInline{ACK=120}。
\item \textbf{情景三：累积确认掩盖前序 ACK 丢失（Cumulative ACK）}。同上两段连续发出后，第一个确认 \StudyInline{ACK=100} 在途中丢失，但第二个确认 \StudyInline{ACK=120} 在第一段超时之前到达了 A。A 读到 $120 > \mathrm{SendBase}$ 即知字节 $92..119$ 全部安全送达，直接将 \StudyInline{SendBase} 推进到 $120$，从而避免了不必要的重传。
\item \textbf{情景四：三个冗余 ACK 触发快速重传（Fast Retransmit）}。超时间隔往往设为百毫秒量级，若流水线中间某一段（如 \StudyInline{Seq=100}）丢失，而后面紧跟的四段（如 \StudyInline{Seq=120, 140, 160, 180}）陆续到达 B，B 每收到一段乱序高序号数据，都不能推进期待前沿，只能重复发回同一个期望号 \StudyInline{ACK=100}。TCP 规定：发送方先收到原本的 1 个正常 \StudyInline{ACK=100}，随后只要再收到 \textbf{3 个冗余确认（triple duplicate ACKs，前后合计收到 4 次 \StudyInline{ACK=100}）}，就不必苦等定时器到期，立即重传缺失的 \StudyInline{Seq=100} 报文段。
\end{itemize}

\textbf{独立变式 / Check：}主机 A 已确认收妥至字节 $499$，随后连续向 B 发送三段数据，长度依次为 $100\text{ B}$、$100\text{ B}$、$150\text{ B}$。若第 1 段途中丢失，第 2、3 段按序到达 B，写出三段的起始序号以及 B 针对第 2、3 段分别发回的确认号。 / Host A has received ACKs through byte 499 and sends three back-to-back segments of 100, 100, and 150 bytes. Segment 1 is lost while Segments 2 and 3 arrive in order. State the three sequence numbers and the two ACK numbers sent by B.

\begin{StudySolution}{答案 / Answer}
三段起始序号依次为 $\mathrm{Seq}_1 = 500$（字节 $500..599$）、$\mathrm{Seq}_2 = 600$（字节 $600..699$）、$\mathrm{Seq}_3 = 700$（字节 $700..849$）。因第 1 段未到，B 期待的最早缺失字节始终停留在 $500$，故收到第 2 段和第 3 段时均发回冗余确认 $\mathrm{ACK} = 500$。 / Sequence numbers are 500, 600, and 700; B sends duplicate $\mathrm{ACK}=500$ upon receiving Segment 2 and again upon receiving Segment 3.
\end{StudySolution}


\StudyHeading{1}{2. 流量控制（Flow Control）｜根据接收缓存余量调节发送上限}
\StudyTarget{CS5222-Chapter03Part3__flow-control}

可靠重传解决了链路丢包，但还没有照顾接收端主机的消化能力。接收端网卡收妥按序字节后，会先放入该连接专属的 TCP 接收缓存（容量记为 \StudyInline{RcvBuffer}），等待上层应用进程调用 \StudyInline{recv()/read()} 取走。如果应用进程正忙于磁盘读写或界面渲染，迟迟不取数据，而发送端持续高速灌入，接收缓存就会溢出。\textbf{流量控制（flow control）}就是让发送速率匹配接收端应用读取速率的端到端速度匹配服务。

记接收缓存中已按序存入但尚未被应用取走的字节量为 $\mathrm{LastByteRcvd} - \mathrm{LastByteRead}$，则当前剩余空闲空间就是\textbf{接收窗口（receive window，记为 \StudyInline{rwnd}）}：$\mathrm{rwnd} = \mathrm{RcvBuffer} - (\mathrm{LastByteRcvd} - \mathrm{LastByteRead})$。接收方在发回的每个 TCP 报文段首部 16 位窗口字段中实时通告最新的 \StudyInline{rwnd}，发送方只需限制在途未确认字节数不超过该值：
\[
\mathrm{LastByteSent} - \mathrm{LastByteAcked} \le \mathrm{rwnd}.
\tag{3.4}
\]
将接收端能力 \StudyInline{rwnd} 与第 5 节即将介绍的网络侧拥塞窗口 \StudyInline{cwnd} 合在一起，TCP 发送方任意时刻允许留在网上的未确认字节上限为：
\[
\mathrm{LastByteSent} - \mathrm{LastByteAcked} \le \min\{\mathrm{cwnd},\, \mathrm{rwnd}\}.
\tag{3.5}
\]
阅读判断选项时注意三处边界：第一，\StudyInline{rwnd} 随着接收缓存的填入与读出动态涨落，并非连接建立后恒定不变；第二，TCP 首部自带 16 位 \StudyInline{rwnd} 字段；第三，当接收缓存填满通告 $\mathrm{rwnd}=0$ 时，若发送方彻底静默且接收方没有任何数据要发，后续缓存腾空时的非零窗口通知一旦丢失就会死锁，因此 TCP 规定发送方在 $\mathrm{rwnd}=0$ 期间仍定期发送携带 $1$ 字节数据的零窗口探测段（zero-window probe），诱发接收方回传最新 \StudyInline{rwnd}。


\StudyHeading{1}{3. 连接管理｜三次握手建立与四次挥手释放}
\StudyTarget{CS5222-Chapter03Part3__connection-mgmt}

在正式交换应用字节前，TCP 通信双方需要协商各自的初始字节序号（ISN）、分配收发缓存并确认对端存活。课件先说明为何简单的“两次握手”在存在可变时延与重传的网络里不安全：若客户端早期发出的连接请求因超时重传，而在本次会话早已结束后那个迟到的旧请求才抵达服务器，仅靠两次握手会让服务器单方面建立“半连接（half-open connection）”，甚至把旧重传数据当作新会话接收。标准 TCP 采用\textbf{三次握手（three-way handshake）}建连，并采用双向独立的\textbf{四次挥手}拆连：

\StudyImage{CS5222/course-notes/assets/ch3p3-tcp-handshake-teardown.png}
\StudyCaption{原创绘制：左图为 TCP 三次握手建立连接（注意 \StudyInline{SYN} 与 \StudyInline{FIN} 标志各消耗 1 个逻辑字节序号，且第 3 步报文段即可捎带首批应用层数据）；右图为四次挥手拆除连接及主动关闭端的 \StudyInline{TIME\_WAIT} 等待。}

\begin{itemize}
\item \textbf{三次握手过程}：（1）客户端选择随机初始序号 $x$，发送纯控制段 \StudyInline{SYN=1, Seq=x}（不携应用数据），进入 \StudyInline{SYN\_SENT}；（2）处于 \StudyInline{LISTEN} 的服务器收到后分配连接缓存与变量，选择自身初始序号 $y$，回复 \StudyInline{SYN=1, ACK=1, Seq=y, ACKnum=x+1}（称 \StudyInline{SYNACK}，同样不携应用数据），进入 \StudyInline{SYN\_RCVD}；（3）客户端收到后分配缓存、进入 \StudyInline{ESTAB}，并回发 \StudyInline{SYN=0, ACK=1, Seq=x+1, ACKnum=y+1}——此时 \StudyInline{SYN} 置 0，已允许在第三步报文段内直接捎带 HTTP \StudyInline{GET} 等应用层请求数据；服务器收到后亦转入 \StudyInline{ESTAB}。
\item \textbf{四次挥手与 \StudyInline{TIME\_WAIT} 的两个作用}：TCP 是全双工通道，任一端完成发送后发 \StudyInline{FIN=1} 只关闭本方向数据流。主动关闭方收到对方的 \StudyInline{FIN} 并发出最后一个 \StudyInline{ACK} 后，不能立刻销毁套接字，而要进入 \StudyInline{TIME\_WAIT} 状态等待 $2\times\mathrm{MSL}$（两倍最大报文段生存时间）：一是防止最后那个 \StudyInline{ACK} 丢失时对端重传 \StudyInline{FIN} 无人应答；二是让本次连接残留在线路中的迟到报文段全部自然过期，避免混入随后复用相同四元组的新连接。
\end{itemize}


\StudyHeading{1}{4. 拥塞控制基本原理｜拥塞的三种代价与两类反馈架构}
\StudyTarget{CS5222-Chapter03Part3__congestion-principles}

初学者常把流量控制与拥塞控制混为一谈。区分两者的标准只看“受保护的对象是谁”：流量控制面向\textbf{点对点的单一接收端主机}，靠接收方明确回告的 \StudyInline{rwnd} 工作；\textbf{拥塞控制（congestion control）}面向\textbf{网络内部多条流共享的路由器队列与链路容量}（即过多源端以过高速率注入数据，超出网络转发能力）。网络拥塞表现为两个典型信号：分组在路由器输出缓存排队造成\textbf{长时延}，以及缓存满溢造成\textbf{分组丢弃}。

课件第 21–29 张幻灯片用三个理想化到现实化的递进模型，算清了为什么放任拥塞不仅会拉大延迟，还会损毁有效吞吐率：
\begin{enumerate}
\item \textbf{场景 1（两条流共享容量 $R$、交换机缓存无限大、无重传）}：每条流的已交付吞吐上限虽能逼近 $R/2$，但当发送速率 $\lambda_{\mathrm{in}}$ 逼近 $R/2$ 时，平均排队长度与端到端时延发散至无穷大。
\item \textbf{场景 2（有限缓存导致丢包 + 发送端重传）}：区分应用层原始注入速率 $\lambda_{\mathrm{in}}$ 与算上重传后的链路供给负载 $\lambda'_{\mathrm{in}} \ge \lambda_{\mathrm{in}}$。只要出现因丢包触发的重传，甚至因排队过长引发\textbf{过早超时重传尚未丢失的包}，宝贵的瓶颈容量 $R$ 就有一部分被用来反复搬运重复副本，导致真正交付给接收端的新数据吞吐率 $\lambda_{\mathrm{out}}$ 显著低于 $R/2$。
\item \textbf{场景 3（多跳路由器 + 有限缓存）}：分组穿过上游多跳路由器时已经消耗了上游链路的转发带宽；一旦它在下游拥塞路由器因队列满而被丢弃，\textbf{此前所有上游跳段为它付出的传输容量全部化为乌有}。当各流拼命加码发送（$\lambda'_{\mathrm{in}}\to\infty$）时，下游队列几乎被注定半途夭折的包占满，全网有效吞吐率 $\lambda_{\mathrm{out}} \to 0$，形成拥塞崩溃（congestion collapse）。
\end{enumerate}

抑制拥塞有两条技术路线：一是传统 TCP 采用的\textbf{端到端拥塞控制（end-to-end）}，网络层不提供显式信令，发送端仅通过观察超时或 3 个冗余 ACK 等丢包现象自行推断拥塞；二是现代数据中心与互联网日益普及的\textbf{网络辅助拥塞控制（network-assisted，如 TCP ECN）}，拥塞路由器在转发 IP 数据报时将首部两位 ECN 字段置为拥塞经历标记，接收端收到后在返回 ACK 首部置位 \StudyInline{ECE} 通知发送端提前缩窗，无需等到缓存溢出丢包。


\StudyHeading{1}{5. TCP 拥塞控制：AIMD、Tahoe 与 Reno、稳态吞吐率及公平性}
\StudyTarget{CS5222-Chapter03Part3__tcp-cc}

TCP 发送端在本地维护拥塞窗口 \StudyInline{cwnd} 与慢启动阈值 \StudyInline{ssthresh}。忽略接收窗口限制时，一个往返时延（RTT）内大约可发出 $\mathrm{cwnd}$ 字节，对应平均发送速率约为 $\mathrm{cwnd}/\mathrm{RTT}$。TCP 通过“未丢包时逐步试探加窗、一见丢包立即缩窗”的闭环自适应调节速率，包含三个状态阶段：

\begingroup\small\setlength{\tabcolsep}{3pt}\renewcommand{\arraystretch}{1.22}
\begin{longtable}{@{}>{\raggedright\arraybackslash}p{\dimexpr0.20\linewidth-1.50000\tabcolsep\relax}>{\raggedright\arraybackslash}p{\dimexpr0.22\linewidth-1.50000\tabcolsep\relax}>{\raggedright\arraybackslash}p{\dimexpr0.28\linewidth-1.50000\tabcolsep\relax}>{\raggedright\arraybackslash}p{\dimexpr0.30\linewidth-1.50000\tabcolsep\relax}@{}}
\toprule
状态阶段 & 进入与维持条件 & 窗口增长规律 & 遇到丢包事件后的状态转移 \\
\midrule\endfirsthead
状态阶段 & 进入与维持条件 & 窗口增长规律 & 遇到丢包事件后的状态转移 \\
\midrule\endhead
\bottomrule\endfoot
\textbf{慢启动（Slow Start, SS）} & 连接刚建立或发生 Timeout 后，且 $\mathrm{cwnd} < \mathrm{ssthresh}$ & \textbf{每 RTT 指数翻倍}：从 $1\text{ MSS}$ 起步，每收到 1 个确认新数据的 ACK 令 $\mathrm{cwnd} \mathrel{+}= 1\text{ MSS}$（即 $1 \to 2 \to 4 \to 8\dots$） & 若涨至 $\mathrm{ssthresh}$ 转入 CA；若发生 Timeout：置 $\mathrm{ssthresh}=\lfloor \mathrm{cwnd}/2\rfloor,\; \mathrm{cwnd}=1$；若遇 3 dup ACKs：Reno 转入 FR，Tahoe 置 1 \\
\textbf{拥塞避免（Congestion Avoidance, CA）} & $\mathrm{cwnd} \ge \mathrm{ssthresh}$ & \textbf{每 RTT 线性加 1（加性增 AI）}：每收到 1 个新 ACK 令 $\mathrm{cwnd} \mathrel{+}= \mathrm{MSS}\cdot(\mathrm{MSS}/\mathrm{cwnd})$，整轮约增加 $1\text{ MSS}$ & 发生 Timeout：置 $\mathrm{ssthresh}=\lfloor \mathrm{cwnd}/2\rfloor,\; \mathrm{cwnd}=1$ 重进 SS；遇 3 dup ACKs：Reno 转入 FR，Tahoe 置 1 重进 SS \\
\textbf{快速恢复（Fast Recovery, FR，仅 Reno）} & 收到 \textbf{3 个冗余 ACK}（后续包仍能到达，表明仅轻度拥塞） & 置 $\mathrm{ssthresh}=\lfloor \mathrm{cwnd}/2\rfloor$，令 \textbf{$\mathrm{cwnd} = \mathrm{ssthresh} + 3\text{ MSS}$}（计入已离开网络的 3 个乱序段）；按课堂与 Tutorial 7 离散轮次口径，随后每轮线性 $+1\text{ MSS}$ & 若恢复期间发生 Timeout，置 $\mathrm{ssthresh}=\lfloor \mathrm{cwnd}/2\rfloor,\; \mathrm{cwnd}=1$ 退回 SS \\
\end{longtable}\endgroup

比较早期版本 \textbf{TCP Tahoe} 与现代基准 \textbf{TCP Reno}，唯一分水岭就在面对“3 个冗余 ACK”时的动作：超时（Timeout）意味网络可能严重堵塞，两者都把阈值减半并将 $\mathrm{cwnd}$ 降回 $1\text{ MSS}$ 重新慢启动；而收到 3 个冗余 ACK 证明管道仍在持续传包，Tahoe 仍一律降回 $1\text{ MSS}$，Reno 则采用\textbf{乘性减（Multiplicative Decrease）}配合快速恢复，把窗口维持在减半阈值加 $3\text{ MSS}$ 的水平继续线性增长，避免了排空管道带来的吞吐骤降。

\StudyImage{CS5222/course-notes/assets/ch3p3-tcp-reno-cwnd.png}
\StudyCaption{原创绘制（对应 Tutorial 7 P1 的 26 轮完整轨迹）：对比 TCP Reno（深蓝实线，第 16 轮在 $\mathrm{cwnd}=42$ 遇 3 个冗余 ACK 后跳至 $21+3=24$ 进入线性增长，第 22 轮在 $\mathrm{cwnd}=29$ 遇超时后置 $\mathrm{ssthresh}=14,\; \mathrm{cwnd}=1$）与 TCP Tahoe（橙色虚线，第 16 轮遇 3 个冗余 ACK 亦降为 1，经慢启动在第 22 轮到达 $\mathrm{ssthresh}=21$ 后转为每轮 $+1$ 线性增长）。}

\textbf{锯齿稳态平均吞吐率与单位换算：}忽略短暂的慢启动阶段，设单条长连接在拥塞避免期呈周期性 AIMD 锯齿波动——窗口每次从 $W/2$（字节）线性爬升至峰值 $W$（字节）触发 3 个冗余 ACK，随即腰斩回 $W/2$。由于线性爬升的平均高度恰好是上下界均值 $\frac{1}{2}(W/2 + W) = \frac{3}{4}W$，稳态平均吞吐率为：
\[
\mathrm{Average\ TCP\ Throughput} = \frac{3}{4}\,\frac{W}{\mathrm{RTT}} \quad (\text{bytes/s}).
\tag{3.6}
\]
做计算题时先看清题目给出的峰值窗口是按段数（MSS）还是按字节给出：若给定峰值窗口为 $W_{\mathrm{seg}} = 64\text{ MSS}$、每段最大净荷 $\mathrm{MSS} = 1460\text{ Bytes}$、$\mathrm{RTT} = 100\text{ ms} = 0.1\text{ s}$，先换算字节窗口 $W = 64\times 1460 = 93{,}440\text{ Bytes}$，再乘以 $8\text{ bits/Byte}$ 转成十进制比特率（网络速率统一采用 $1\text{ Mbps} = 10^6\text{ bps},\; 1\text{ kbps} = 10^3\text{ bps}$）：
\[
\mathrm{Average\ Throughput} = \frac{3}{4}\times \frac{64\times 1460\times 8\text{ bits}}{0.1\text{ s}} = 5{,}606{,}400\text{ bps} = 5.6064\text{ Mbps}.
\tag{3.7}
\]

\textbf{AIMD 的几何公平性与两项现实限制：}当两条具有相同 RTT 的 TCP 连接共享容量为 $R$ 的瓶颈链路时，在 $(R_1, R_2)$ 二维吞吐平面上，加性增让工作点沿斜率为 $1$ 的 $45^\circ$ 方向右移上移（双方等量增加），越过容量线 $R_1 + R_2 = R$ 丢包后，乘性减让工作点沿连向原点的射线折半回缩（高吞吐一方被削减得更多）。反复“平行 $45^\circ$ 推进 $\to$ 向原点折半”，状态点必然收敛到等分线 $R_1 = R_2 = R/2$ 上来回振荡。但这种公平保证存在两个常考例外（详见\StudyRef{CS5222-Tutorial07}{p2}{Tutorial 7 P2}）：一是流媒体若采用无拥塞控制的 \textbf{UDP} 恒速发包，拥塞时不会主动退让，可能挤压 TCP 带宽；二是公平性仅约束\textbf{单条 TCP 连接}而非单个应用程序——若瓶颈上已有 $9$ 条连接，某个浏览器新建 $1$ 条连接可得 $R/10$，若并发开启 $11$ 条连接则可分走 $11R/20 > R/2$。

\textbf{独立变式 / Transfer：}某 TCP Reno 连接当前处于拥塞避免阶段，$\mathrm{cwnd} = 30\text{ MSS}$，$\mathrm{RTT} = 50\text{ ms}$，$1\text{ MSS} = 1500\text{ Bytes}$。（1）求该窗口在 $15\text{ MSS} \leftrightarrow 30\text{ MSS}$ 间稳态锯齿波动时的平均吞吐率（以 Mbps 计）；（2）若在 $\mathrm{cwnd}=30\text{ MSS}$ 时收到 3 个冗余 ACK，求更新后的 $\mathrm{ssthresh}$ 与 $\mathrm{cwnd}$。 / A TCP Reno flow oscillates between 15 and 30 MSS with $\mathrm{RTT}=50\text{ ms}$ and $\mathrm{MSS}=1500\text{ B}$. Find its steady-state average throughput in Mbps, and give the new $\mathrm{ssthresh}$ and $\mathrm{cwnd}$ upon triple duplicate ACKs at 30 MSS.

\begin{StudySolution}{答案 / Answer}
（1）峰值字节数 $W = 30\times 1500 = 45{,}000\text{ B}$，代入式（3.6）得平均吞吐率 $= \frac{3}{4}\times \frac{45{,}000\times 8}{0.05} = 5.4\times 10^6\text{ bps} = 5.4\text{ Mbps}$。（2）收到 3 个冗余 ACK 时，Reno 置 $\mathrm{ssthresh} = \lfloor 30/2\rfloor = 15\text{ MSS}$，$\mathrm{cwnd} = 15 + 3 = 18\text{ MSS}$。 / Average throughput is $\frac{3}{4}\frac{30\times 1500\times 8}{0.05} = 5.4\text{ Mbps}$; on triple duplicate ACKs, $\mathrm{ssthresh}=15\text{ MSS}$ and $\mathrm{cwnd}=18\text{ MSS}$.
\end{StudySolution}


\StudyHeading{1}{6. 课堂 Q\&A 完整原题与双语解析（Chapter 3 全 7 题）}
\StudyTarget{CS5222-Chapter03Part3__qa}

本节收录 \texttt{Q\&A\_cha3\_2026\_student.pptx}（第 2–8 张幻灯片）全部 7 道课堂选择题的完整英文题干与选项（前 5 题回扣\StudyRef{CS5222-Chapter03}{rdt}{Part 1–2 的 rdt 与 GBN/SR}，后 2 题考查本篇 TCP Reno 窗口跳变）：

\begin{enumerate}
\item \textbf{Q1 (Slide 2):} \emph{Suppose a sender sends a packet to the receiver. Which of the following mechanism can a \textbf{receiver} use to determine whether the packet has been \textbf{corrupted}? (Choose one)} \\
\texttt{A) ACK \quad B) Checksum \quad C) Timeout \quad D) None of the above} \\
\textbf{解析 / Answer: B) Checksum。}接收方对收到的分组重新计算检验和并与首部字段比对，检验比特翻转差错。 / The receiver verifies bit-level integrity using the packet's checksum.

\item \textbf{Q2 (Slide 3):} \emph{Suppose a sender sends a packet to the receiver. Which of the following mechanism can a \textbf{sender} use to determine whether the packet has been \textbf{corrupted}? (Choose one)} \\
\texttt{A) ACK \quad B) Checksum \quad C) Timeout \quad D) None of the above} \\
\textbf{解析 / Answer: A) ACK。}数据包离开后发送方自己看不见接收端波形，只能依靠接收方发回的确认反馈（如 \texttt{rdt2.2/TCP} 中的冗余旧 ACK 或显式 NAK）间接获知前向分组受损或失序。 / The sender relies on receiver feedback (duplicate ACKs or NAKs) to learn about forward-path corruption.

\item \textbf{Q3 (Slide 4):} \emph{Suppose a sender sends a packet to the receiver. Which of the following mechanism can a \textbf{sender} use to determine \textbf{for sure} whether the packet has been \textbf{lost}? (Choose one)} \\
\texttt{A) ACK \quad B) Checksum \quad C) Timeout \quad D) None of the above} \\
\textbf{解析 / Answer: D) None of the above。}注意题干限定词 \emph{for sure}（百分之百确定）：即使定时器超时，也可能是数据包仍在异常缓慢的队列中传输，或者数据包早已安全抵达而仅仅是回程的 ACK 丢了，发送方无法仅凭本地超时断定原包必然丢失。 / Neither timeout nor missing ACKs proves packet loss with certainty, since extreme delay or a lost ACK produces the exact same symptom at the sender.

\item \textbf{Q4 (Slide 5):} \emph{Which of the following statement(s) about \textbf{GBN} is/are NOT true?} \\
A) It is a pipelined protocol. \\
B) The sender needs one timer for its operation. \\
C) The receiver needs one buffer to hold some out-of-order packet. \\
D) When timeout, the sender just retransmits that unacked packet causing the timeout. \\
\textbf{解析 / Answer: C and D are NOT true。}GBN 接收端只维护单个期望序号变量 \StudyInline{expectedseqnum}，\textbf{不缓存任何乱序到达的分组}（直接丢弃并重发旧累积 ACK，故 C 错）；超时发生时，GBN 发送端重传从最早未确认分组起的\textbf{所有已发未确认分组}（故 D 错，只重传单个超时分组的是 SR 或简化 TCP）。 / C and D are false: a GBN receiver discards out-of-order packets without buffering, and a GBN sender retransmits all unacknowledged packets in the window upon timeout.

\item \textbf{Q5 (Slide 6):} \emph{Which of the following statement(s) about \textbf{SR} is/are NOT true?} \\
A) It is a pipelined protocol. \\
B) The sender needs one timer for its operation. \\
C) The receiver needs one buffer to hold some out-of-order packet. \\
D) When timeout, the sender just retransmits that unacked packet causing the timeout. \\
\textbf{解析 / Answer: B is NOT true。}选择重传（SR）只重传各自超时的独立分组，因此在概念模型上发送端需要为\textbf{每一个在途未确认分组分别维护独立的逻辑定时器}，只有一个总定时器不够（注意选项 C 中的 ``one buffer'' 指接收端设有一块接收缓存区用来暂存乱序包，符合 SR 原理）。 / B is false: Selective Repeat maintains a separate logical timer per unacknowledged packet rather than a single shared timer.

\item \textbf{Q6 (Slide 7):} \emph{Consider \textbf{TCP Reno}. Suppose the sender observes \textbf{three duplicated ACKs} at around 15. What will be the congestion window size at around 16?} \\
\texttt{A) 1 \quad B) 3 \quad C) 7 \quad D) 10} \\
\textbf{解析 / Answer: D) 10。}原幻灯片中的 ``at around 15'' 对应课堂拥塞控制示例图第 15 传输轮次（”at round 15“ 的口语笔误，该图第 15 轮时 $\mathrm{cwnd} = 14\text{ MSS}$）。收到 3 个冗余 ACK 后，TCP Reno 设 $\mathrm{ssthresh} = 14 / 2 = 7$，第 16 轮窗口变为 $\mathrm{cwnd} = \mathrm{ssthresh} + 3 = 7 + 3 = 10\text{ MSS}$。 / Referring to the lecture cwnd trace where $\mathrm{cwnd}=14$ at round 15, TCP Reno sets $\mathrm{ssthresh}=7$ and $\mathrm{cwnd}=7+3=10$ at round 16.

\item \textbf{Q7 (Slide 8):} \emph{Consider \textbf{TCP Reno}. Suppose the sender experiences \textbf{a timeout} at around 15. What will be the congestion window size at around 16?} \\
\texttt{A) 1 \quad B) 3 \quad C) 7 \quad D) 10} \\
\textbf{解析 / Answer: A) 1。}只要发生定时器超时，无论是 TCP Tahoe 还是 TCP Reno，下一轮都一律将 $\mathrm{cwnd}$ 重置为 $1\text{ MSS}$ 并重新进入慢启动（同时将 $\mathrm{ssthresh}$ 更新为 $7$）。 / Upon a timeout, TCP Reno resets $\mathrm{cwnd}$ to $1\text{ MSS}$ and enters slow start.
\end{enumerate}

\textbf{English takeaway:} Flow control prevents overflowing the receiver buffer via \StudyInline{rwnd}, whereas congestion control prevents overloading network routers via \StudyInline{cwnd}. Under TCP Reno, triple duplicate ACKs trigger fast retransmit and fast recovery ($\mathrm{ssthresh}=\lfloor \mathrm{cwnd}/2\rfloor$, $\mathrm{cwnd}=\mathrm{ssthresh}+3$), whereas a timeout resets \StudyInline{cwnd} to $1\text{ MSS}$ and restarts slow start.


\StudyHeading{1}{学习要求与来源覆盖索引}
\StudyTarget{CS5222-Chapter03Part3__sources}

\begingroup\small\setlength{\tabcolsep}{3pt}\renewcommand{\arraystretch}{1.22}
\begin{longtable}{@{}>{\raggedright\arraybackslash}p{\dimexpr0.30\linewidth-1.00000\tabcolsep\relax}>{\raggedright\arraybackslash}p{\dimexpr0.70\linewidth-1.00000\tabcolsep\relax}@{}}
\toprule
原课件页段（Part 3 / 整章页脚） & 讲义对应位置与核心内容 \\
\midrule\endfirsthead
原课件页段（Part 3 / 整章页脚） & 讲义对应位置与核心内容 \\
\midrule\endhead
\bottomrule\endfoot
第 1–9 页（总 71–79 页） & §1：简化 TCP 发送端三事件、四类重传时序与 3 个冗余 ACK 快速重传 \\
第 10–14 页（总 80–84 页） & §2：接收缓存余量 \texttt{rwnd}、发送窗口双约束与零窗口探测 \\
第 15–20 页（总 85–90 页） & §3：两次握手缺陷、三次握手状态变迁、四次挥手与 \texttt{TIME\_WAIT} \\
第 21–31 页（总 91–101 页） & §4：拥塞三大代价场景（排队发散、重传浪费、多跳上游作废）与端到端/ECN 反馈 \\
第 32–48 页（总 102–118 页） & §5：慢启动、拥塞避免、快速恢复（Tahoe vs Reno）、$\frac{3}{4}W/\mathrm{RTT}$ 吞吐率与 AIMD 公平性 \\
\texttt{Q\&A\_cha3} 第 1–8 页 & §6：检验和、确认反馈、GBN/SR 定时器与缓存辨析、Reno 第 15–16 轮窗口跳变 \\
\end{longtable}\endgroup
